% ===== §7 =====
\section{The Unsayable Direction of Growth}
\label{sec:evolution}

\noindent
This section takes the sixth and last dimension, the evolution of a relation, and establishes that the direction of a bond's genuine growth cannot be stated in the signifiers now available, because genuine novelty is by definition what the present vocabulary does not contain, so that a relation is a becoming undergone rather than a plan carried out. The section proceeds through the phenomenon, a review of the accounts that would make the direction of growth sayable in advance, the isolation of the core, and the claim, and it closes Part One by gathering the six limits into their shared form.

\subsection*{The phenomenon}

A relation grows and turns. It generates forms it did not have, becomes over years something neither party could have set out at the start. Language seems able to reach this growth, at least in prospect. We plan, we imagine, we say where we are heading: what we will build, who we will become, what our life will grow into. These sentences present the direction of a bond's development as if it could be laid out ahead of the walking.

They lay out only what the present already contains.

\subsection*{Two ways of saying the future, and their limit}

Two accounts promise that the direction of growth can be stated in advance, and each shows, by the form of its shortfall, what genuine novelty is.

The first is prediction, the projection of a future from the trends of the present. Prediction is powerful and often right, and it has a definite scope: it states the future insofar as the future is a continuation of the present, an extrapolation of what is already underway. But precisely for that reason prediction reaches only the future that the present already implies. What it cannot reach is the genuinely new, the form that is not an extrapolation of present trends but a break from them, for such a form is not contained in the present data and so cannot be projected from them. Prediction states the future that is already latent; the direction of genuine growth is the future that is not latent, and prediction passes it by.

The second is the account of the future as the realization of a plan, the execution of an intended design. A couple sets out to build a life of a certain shape, and the building realizes the design. This too has its truth, and this too has its limit. A plan can only specify what its authors can already conceive, in the vocabulary already available to them; it lays out a future in present terms. But a relation that genuinely evolves becomes something its members could not have specified at the outset, something for which they did not yet have the words, and this is not a failure of planning but the mark of real growth. Where the outcome was fully specifiable in advance, there was no evolution, only the unrolling of a design. The plan reaches the future it could already state; the direction of genuine growth is the future it could not, and the plan does not contain it.

\subsection*{The core}

The core language cannot reach here is the \emph{direction} of the relation's becoming, its genuine novelty, which cannot be stated in the signifiers now available. The point is not that the future is uncertain in the ordinary way, a known range of outcomes not yet decided; uncertainty of that kind is still stated in present terms, as a distribution over foreseen possibilities. The point is that the genuinely new form is not yet sayable at all, since its sayability would require the very signifiers whose absence is what makes it new. The direction of real growth is the not-yet-symbolizable, and it stays ahead of speech not by accident but by what novelty is. This is the forward-facing form of the same generativity that justice was shown to guard: there, the still-forming as the object of protection; here, the not-yet-sayable as the direction of growth.

\begin{claim}[The direction of growth is not sayable in advance]
The direction in which a relation genuinely evolves cannot be stated in the signifiers now available, because genuine novelty is precisely what the present vocabulary does not contain. Prediction reaches only the latent future, the plan only the already-conceivable; the new form is not-yet-symbolizable, and stays ahead of every advance description by what novelty is.
\end{claim}

\begin{casebox}{the life neither could have described}
Two people set out with a picture of the life they would build, and years on they are something the picture did not contain, not a deviation from the plan but a growth the plan had no words for. Asked at the start to describe what they became, they could not have, because the terms were not yet theirs. The becoming was real precisely in that it outran the description available to it. What they could not have said in advance was not a gap in their foresight. It was the room in which the relation grew into what no one yet had words for.
\end{casebox}

\subsection*{Why the limit is generative here}

This is what keeps the evolution of a relation open rather than executed. Because the direction of genuine growth cannot be said in advance, a relation is not a plan carried out but a becoming undergone, and its future is generated in the living rather than read from a design. Consider the counter-case, a relation whose development could be fully stated beforehand. It would have no future in the strong sense, only the unrolling of a foreseen sequence, a life lived as the execution of its own prospectus. The limit forbids this. Because the direction of real growth is not-yet-sayable, the future of a relation cannot be pre-stated and so cannot be merely executed; it must be undergone, generated in the walking, open to forms not yet in the language. What cannot be said in advance is not the failure of foresight. It is the space in which a relation can still become what no one yet has words for.

\begin{figure}[t!]
\centering
\begin{tikzpicture}[scale=1.0,
  dim/.style={draw=rosedeep,thick,rounded corners=2pt,fill=pageblush,inner sep=4pt,align=center,font=\footnotesize\color{inksoft},text width=2.35cm},
  core/.style={draw=gold,very thick,circle,fill=pageblush,inner sep=2pt,align=center,font=\footnotesize\itshape\color{rosedeep},text width=1.5cm},
  reach/.style={-{Latex[length=2mm]},rose,thick}]
\node[core] (m) at (0,0) {manque};
\foreach \ang/\lab in {90/{Existence\\ \tiny the being of the bond},
                       30/{Maintenance\\ \tiny the living present},
                       -30/{Value\\ \tiny no place to price it},
                       -90/{Ethics\\ \tiny the aesthetic measure},
                       -150/{Justice\\ \tiny generativity guarded},
                       150/{Evolution\\ \tiny the unsayable direction}}{
  \node[dim] (\ang) at (\ang:3.5cm) {\lab};
  \draw[reach] (\ang) -- ($(m)!0.62!(\ang)$);
}
\end{tikzpicture}
\caption{Six windows on one gap. In each dimension of relational life, language reaches toward a core it does not touch (rose arrows), and in each the core is a face of the same manque. The parallel is not a ranking; it is six structured phenomena sharing one form.}
\label{fig:sixwindows}
\end{figure}

\goldrule
\noindent
The six limits are now before us, each with its own core and each with its own claim: the ungroundability of the bond's existence, the uncapturability of the living present, the absence of any external place to price relational value, the aesthetic measure of the ethical response, the generativity that justice guards, and the unsayable-in-advance direction of growth. They are set side by side, not ranked and not derived one from another, and the parallel is deliberate: the limit of language is not one vague insufficiency but six distinct and structured phenomena. Yet across all six a single form has appeared. In each, language reaches toward something and does not reach it; in each, the core it cannot reach is a remainder that the operation of the signifier leaves outside itself; and in each, the not-reaching is not a defect but the condition under which that dimension of relational life is what it is. Part Two asks what stands at these limits, traces it to the constitution of the speaking subject, and gives it a name.
